% =============================================================================
% Kapitel 7: Fazit und Ausblick
% =============================================================================
\chapter{Fazit und Ausblick}
\label{ch:fazit}

\section{Zusammenfassung}

In dieser Arbeit wurde die Software \evax{} eingesetzt, um die lokale
atomare Struktur zweier Probensysteme aus EXAFS-Daten zu bestimmen:
einer ternären CoNiCu-Mittelentropielegierung und bimetallischer
CoPt-Nanopartikel.

Die methodische Vorgehensweise umfasste die sorgfältige Extraktion
der $\chi(k)$-Daten aus Multi-Edge-Absorptionsspektren, den Aufbau
geeigneter Startstrukturen, eine systematische Parameterstudie zur
Identifikation optimaler Fit-Einstellungen sowie den Vergleich
verschiedener Kristallstrukturen (FCC, BCC, HCP).

\TODO{Nach Abschluss der Fits: Kernaussagen zu den Ergebnissen
zusammenfassen, z.\,B.:}
\begin{itemize}[nosep]
  \item \TODO{Die CoNiCu-MEA kristallisiert in der FCC-Struktur
        ($R_{\text{total}} = \ldots$ vs.\ $\ldots$ für BCC).}
  \item \TODO{Die chemische Nahordnung zeigt eine leichte
        Bevorzugung von Co-Ni-Paaren ($\alpha \approx -0{,}10$),
        während Co-Co-Paare vermieden werden
        ($\alpha \approx +0{,}11$).}
  \item \TODO{Die CoPt-Nanopartikel weisen \ldots}
\end{itemize}

Die Parameterstudie belegt, dass der Wavelet-Raum als Vergleichsraum,
eine $k^2$-Gewichtung und die Berücksichtigung von
Mehrfachstreupfaden die besten Ergebnisse liefern.
Die strukturellen Befunde erwiesen sich als robust gegenüber
Variationen im $k$-Bereich und der Schrittweite.

\section{Bewertung der Methode}

\evax{} hat sich als geeignetes Werkzeug für die Analyse komplexer
Mehrstoffsysteme erwiesen.
Die simultane Multi-Edge-Auswertung überwindet die bei konventionellen
Fits problematische Parameterkorrelation und ermöglicht die direkte
Bestimmung elementaufgelöster Strukturparameter.

Die Haupteinschränkung liegt in der Rechenzeit:
Ein einzelner Produktionslauf erfordert mehrere Stunden, und die
systematische Parameterexploration vervielfacht diesen Aufwand.
Die in dieser Arbeit entwickelte Parallelisierungsstrategie --
isolierte Arbeitsverzeichnisse mit automatisierter Konfiguration --
hat sich als praktikable Lösung bewährt und ließ sich auf einem
handelsüblichen Laptop mit sechs gleichzeitigen Prozessen umsetzen.

\section{Ausblick}

Mehrere Erweiterungen bieten sich für zukünftige Arbeiten an:

\begin{itemize}
  \item \textbf{Temperaturabhängige Messungen:}
    Durch EXAFS-Messungen bei mehreren Temperaturen ließen sich
    thermische und statische Beiträge zum Debye-Waller-Faktor
    trennen und die Debye-Temperatur elementaufgelöst bestimmen.
  \item \textbf{Finite-Size-Modelle:}
    Für die CoPt-Nanopartikel könnte ein endlicher Cluster statt
    einer periodischen Superzelle verwendet werden, um
    Oberflächeneffekte korrekt zu erfassen.
  \item \textbf{Kombination mit XRD:}
    Die Gitterkonstante aus der Röntgendiffraktometrie kann als
    zusätzlicher Constraint in die \evax{}-Optimierung einfließen
    und die Mehrdeutigkeit weiter reduzieren.
  \item \textbf{Vergleich mit DFT-Rechnungen:}
    Ab-initio-Rechnungen der Mischungsenergie und
    SRO-Gleichgewichtsparameter könnten die experimentell bestimmte
    Nahordnung unabhängig validieren.
\end{itemize}
